\section{Kadison--Schwarz inequality}

A linear map is completely positive (Definition~\ref{def:cp_map}) if and only
if it has a Kraus form. The Kadison--Schwarz inequality and the
multiplicative-domain characterisation are proved first for Kraus maps.
The transfer map $\E_A(X) = \sum_i A^i X (A^i)^\dagger$ is a Kraus map
(Theorem~\ref{thm:transfer_cp}), so these results apply directly to transfer
maps.

\begin{definition}[Kraus map]\label{def:kraus_map}
    \lean{KadisonSchwarz.krausMap}
    \leanok
    \uses{def:cp_map}
    Given operators $\{K_i\}_{i=0}^{d-1}$ with $K_i \in \MN{D}$,
    the \emph{Kraus map} is
    $E(X) = \sum_{i=0}^{d-1} K_i X K_i^\dagger$.
    A linear map is completely positive
    (Definition~\ref{def:cp_map}) if and only if it can be written
    in this form.
\end{definition}

\begin{definition}[Adjoint Kraus map]\label{def:kraus_adjoint_map}
    \lean{KadisonSchwarz.krausAdjointMap}
    \leanok
    \uses{def:kraus_map}
    The \emph{adjoint Kraus map} is
    $E^*(X) = \sum_{i=0}^{d-1} K_i^\dagger X K_i$.
    When the $\{K_i\}$ are the matrices of an MPS tensor~$A$,
    this is the transfer map of the conjugate-transposed family
    $i \mapsto (A^i)^\dagger$.
\end{definition}

\begin{definition}[Unital Kraus map]\label{def:unital_kraus}
    \lean{KadisonSchwarz.IsUnitalKraus}
    \leanok
    \uses{def:kraus_map}
    A Kraus map is \emph{unital} if
    $\sum_{i=0}^{d-1} K_i K_i^\dagger = \Id$.
    Equivalently, $E(\Id) = \Id$.
\end{definition}

\begin{definition}[Trace-preserving Kraus map]\label{def:tp_kraus}
    \lean{KadisonSchwarz.IsTPKraus}
    \leanok
    \uses{def:kraus_map}
    A Kraus map is \emph{trace-preserving} if
    $\sum_{i=0}^{d-1} K_i^\dagger K_i = \Id$.
    Equivalently, the adjoint Kraus map is unital.
    This is the standard MPS normalization condition.
    In the later gauge language it is the left-canonical condition,
    so Kadison--Schwarz arguments are often applied to the adjoint map.
\end{definition}

\begin{theorem}[Kadison--Schwarz inequality]\label{thm:kadison_schwarz}
    \lean{KadisonSchwarz.kadison_schwarz}
    \leanok
    \uses{def:unital_kraus}
    Let $E(X) = \sum_i K_i X K_i^\dagger$ be a \emph{unital}
    Kraus map (i.e.\ $\sum_i K_i K_i^\dagger = \Id$).
    Then for all $X \in \MN{D}$,
    \[
        E(X^\dagger X) \ge E(X)^\dagger E(X).
    \]
    This is \cite[Equation~(5.2)]{Wolf2012Quantum}.
\end{theorem}

\begin{proof}
    \leanok
    The $2 \times 2$ block matrix
    $[\begin{smallmatrix} X^\dagger X & X^\dagger \\ X & \Id \end{smallmatrix}]$
    is PSD ($= v v^\dagger$ with $v = [X^\dagger,\, \Id]^T$).
    Applying the unital CP map blockwise preserves positive semidefiniteness
    (each block $K_i (\cdot) K_i^\dagger$ preserves PSD, and unitality
    ensures the $(2,2)$-block maps to~$\Id$).
    The Schur complement of the $(2,2)$-block $\Id$ gives the result.
\end{proof}

\begin{theorem}[Kadison--Schwarz for the adjoint channel]\label{thm:kadison_schwarz_adjoint}
    \lean{KadisonSchwarz.kadison_schwarz_adjoint}
    \leanok
    \uses{def:tp_kraus, def:kraus_adjoint_map}
    If $\sum_i K_i^\dagger K_i = \Id$ (trace-preserving), then
    the adjoint Kraus map satisfies
    $E^*(X^\dagger X) \ge E^*(X)^\dagger E^*(X)$.
\end{theorem}

\begin{proof}
    \leanok
    \uses{thm:kadison_schwarz}
    The adjoint operators $\{K_i^\dagger\}$ satisfy
    $\sum_i K_i^\dagger (K_i^\dagger)^\dagger = \sum_i K_i^\dagger K_i = \Id$,
    so they form a unital family.
    Apply Theorem~\ref{thm:kadison_schwarz} to $\{K_i^\dagger\}$.
\end{proof}

\section{Multiplicative domain}

\begin{theorem}[KS gap decomposition]\label{thm:ks_gap_decomp}
    \lean{KadisonSchwarz.ks_gap_eq_sum_squares}
    \leanok
    \uses{def:unital_kraus, def:kraus_map}
    For a unital Kraus map $E$,
    the Kadison--Schwarz gap decomposes as
    \[
        E(X^\dagger X) - E(X)^\dagger E(X)
        =
        \sum_{i=0}^{d-1} R_i^\dagger R_i,
    \]
    where $R_i := X K_i^\dagger - K_i^\dagger E(X)$.
\end{theorem}

\begin{proof}
    \leanok
    \uses{def:unital_kraus}
    Expand the right-hand side
    $\sum_i (X K_i^\dagger - K_i^\dagger E(X))^\dagger
        (X K_i^\dagger - K_i^\dagger E(X))$,
    distribute, and use unitality $\sum_i K_i K_i^\dagger = \Id$
    to identify the result with the KS gap.
\end{proof}

\begin{theorem}[KS equality implies Kraus commutation]\label{thm:kraus_commute}
    \lean{KadisonSchwarz.kraus_commute_of_ks_equality}
    \leanok
    \uses{def:unital_kraus}
    Let $E$ be a unital Kraus map.
    If $E(X^\dagger X) = E(X)^\dagger E(X)$ for some $X$,
    then $X K_i^\dagger = K_i^\dagger E(X)$ for all~$i$.
\end{theorem}

\begin{proof}\leanok\uses{thm:ks_gap_decomp}
    By Theorem~\ref{thm:ks_gap_decomp},
    the KS gap is $\sum_i R_i^\dagger R_i$ with
    $R_i = X K_i^\dagger - K_i^\dagger E(X)$.
    If the gap vanishes, $\sum_i R_i^\dagger R_i = 0$.
    Each summand $R_i^\dagger R_i$ is PSD, so each must be zero,
    giving $R_i = 0$.
\end{proof}

\begin{theorem}[KS equality for peripheral eigenvectors]\label{thm:ks_peripheral}
    \lean{KadisonSchwarz.ks_equality_of_peripheral_eigenvector}
    \leanok
    \uses{def:unital_kraus, def:tp_kraus}
    Let $E$ be a Kraus map that is both unital and trace-preserving.
    If $E(X) = \mu X$ with $|\mu| = 1$, then the KS gap vanishes:
    $E(X^\dagger X) = E(X)^\dagger E(X)$.
\end{theorem}

\begin{proof}
    \leanok
    \uses{thm:kadison_schwarz}
    The gap $G := E(X^\dagger X) - E(X)^\dagger E(X)$ is PSD
    by Theorem~\ref{thm:kadison_schwarz}.
    Its trace satisfies
    $\tr(G) = \tr(E(X^\dagger X)) - \tr(E(X)^\dagger E(X))$.
    Trace preservation gives $\tr(E(X^\dagger X)) = \tr(X^\dagger X)$.
    Using $E(X) = \mu X$ and $|\mu| = 1$,
    $\tr(E(X)^\dagger E(X)) = |\mu|^2 \tr(X^\dagger X) = \tr(X^\dagger X)$.
    So $\tr(G) = 0$, and a PSD matrix with zero trace must be zero.
\end{proof}

\begin{theorem}[Left multiplicative identity from KS equality]\label{thm:left_multiplicative_identity}
    \lean{KadisonSchwarz.multiplicative_domain_left}
    \leanok
    \uses{def:unital_kraus}
    Let $E$ be a unital Kraus map.
    If $E(X^\dagger X) = E(X)^\dagger E(X)$, then
    $E(X^\dagger Y) = E(X)^\dagger E(Y)$ for all $Y \in \MN{D}$.
\end{theorem}

\begin{proof}
    \leanok
    \uses{thm:kraus_commute}
    From Theorem~\ref{thm:kraus_commute},
    $X K_i^\dagger = K_i^\dagger E(X)$ for all~$i$.
    Taking conjugate transposes:
    $K_i X^\dagger = E(X)^\dagger K_i$.
    Then
    \begin{align*}
        E(X^\dagger Y)
         & = \sum_i K_i X^\dagger Y K_i^\dagger    \\
         & = \sum_i E(X)^\dagger K_i Y K_i^\dagger \\
         & = E(X)^\dagger E(Y).
    \end{align*}
\end{proof}

\begin{theorem}[Right multiplicative identity from KS equality]
    \label{thm:right_multiplicative_identity}
    \lean{KadisonSchwarz.multiplicative_domain_right}
    \leanok
    \uses{def:unital_kraus}
    Let $E$ be a unital Kraus map.
    If $E(X^\dagger X) = E(X)^\dagger E(X)$, then
    $E(YX) = E(Y)E(X)$ for all $Y \in \MN{D}$.
\end{theorem}

\begin{proof}
    \leanok
    \uses{thm:kraus_commute}
    From Theorem~\ref{thm:kraus_commute},
    $X K_i^\dagger = K_i^\dagger E(X)$ for all~$i$. Then
    \begin{align*}
        E(YX)
         & = \sum_i K_i Y X K_i^\dagger    \\
         & = \sum_i K_i Y K_i^\dagger E(X) \\
         & = E(Y)E(X).
    \end{align*}
\end{proof}
